\documentclass[11pt]{article}
\usepackage[margin=1in]{geometry}
\usepackage{amsmath,amssymb,amsthm}
\usepackage{booktabs}
\usepackage{hyperref}
\newtheorem{theorem}{Theorem}
\newtheorem{lemma}[theorem]{Lemma}
\newtheorem{proposition}[theorem]{Proposition}
\newtheorem{corollary}[theorem]{Corollary}
\theoremstyle{definition}
\newtheorem{definition}[theorem]{Definition}
\theoremstyle{remark}
\newtheorem{remark}[theorem]{Remark}

\title{Disproof of Conjecture 00000001685:\\
no even subdivision of $K_{3,3}$ is a Pfaffian graph}
\author{jilint777}
\date{2026-10-04}

\begin{document}
\sloppy
\maketitle

\begin{abstract}
Conjecture 00000001685 says that the minimal families of nonplanar Pfaffian graphs
are the M\"obius ladders $M_{2k}$ and the even subdivisions of $K_{3,3}$. In
particular it says that the even subdivisions of $K_{3,3}$ are Pfaffian graphs.
This is false. We prove that \emph{no} even subdivision of $K_{3,3}$ has a Pfaffian
orientation, including $K_{3,3}$ itself, which is also the M\"obius ladder on $6$
vertices. The proof is a parity argument: in every orientation an odd number of
the six perfect-matching terms of the Pfaffian are negative, so the terms never
all have the same sign. The obstruction, and its instances for $K_{3,3}$, for a
proper even subdivision of $K_{3,3}$, and for the $6$-vertex M\"obius ladder, are
formalized in Lean~4 (core library only). As a secondary result we give a
nonplanar Pfaffian graph of crossing number $1$ that lies in neither family, and
we show that the M\"obius clause also fails if the families are read as
non-Pfaffian obstructions.
\end{abstract}

\section{The conjecture and the clause we refute}
\begin{quote}
\emph{Definition:} A Pfaffian orientable graph is one whose perfect matchings can
be counted by signed cancellation. \emph{Conjecture:} The minimal families of
nonplanar Pfaffian graphs are the M\"obius ladders $M_{2k}$ and the even
subdivisions of $K_{3,3}$; these two families exhaust the nonplanar Pfaffian class
of crossing number $\le 2$.
\end{quote}
The Chinese version says the same. The first sentence names two families
\emph{of nonplanar Pfaffian graphs}, and the second calls them part of ``the
nonplanar Pfaffian class''. So the conjecture asserts, at the least:

\begin{quote}
\textbf{Clause (E).} Every even subdivision of $K_{3,3}$ is a (nonplanar) Pfaffian graph.
\end{quote}

The conjecture is a conjunction that includes (E), so refuting (E) refutes the
conjecture. Whatever ``minimal'' means, a family of nonplanar Pfaffian graphs
has to consist of Pfaffian graphs, and the falsity of (E) does not depend on that
word.

\paragraph{Readings covered.}
\begin{itemize}
\item \emph{Even subdivision.} An even subdivision of $H$ replaces every edge by a
path with an even number ($0,2,4,\dots$) of internal vertices. We prove
(Theorem~\ref{thm:main}) that \emph{every} even subdivision of $K_{3,3}$ is
non-Pfaffian. So it makes no difference whether $K_{3,3}$ itself counts as an even
subdivision, whether at least one edge must be subdivided, or whether every
edge must be subdivided.
\item \emph{Pfaffian.} We use the conjecture's own definition (the signed sum
counts the perfect matchings, $|\operatorname{Pf}(A_D)|=\#\mathrm{PM}$) and
Kasteleyn's standard definition (all perfect matchings have the same sign). They
are equivalent (Lemma~\ref{lem:equiv}), and both are refuted.
\item \emph{$M_{2k}$.} If $M_{2k}$ is the M\"obius ladder on $2k$ vertices ($k$
rungs), then $M_6\cong K_{3,3}$ is not Pfaffian, so the M\"obius part fails as well.
More generally, the computations below show that the ladders with $3$, $5$ and $7$
rungs are non-Pfaffian. If instead $M_{2k}$ has $2k$ rungs, the ladders with $2$, $4$
and $6$ rungs are Pfaffian, and we make no claim against the M\"obius part. In
both cases the refutation of (E) does not depend on this choice.
\item \emph{Obstruction reading.} Some readers may take the ``minimal families''
to be minimal \emph{non}-Pfaffian obstructions, in the spirit of Little's
theorem~\cite{little}. Under that reading the M\"obius clause fails in both
conventions. The Wagner graph $V_8$ is Pfaffian (Lean: \texttt{M8\_pfaffian}), and it
lies in the family $M_{2k}$ either way: $k=4$ when counting vertices, and $k=2$ when
counting rungs. Lean: \texttt{obstruction\_reading\_false} (rung count,
$M_{2k}=\texttt{mobius}(2k)$) and
\texttt{obstruction\_reading\_false\_vertex\_convention} ($M_{2k}=\texttt{mobius}\,k$).
\end{itemize}

\section{Definitions}
Let $G$ be a loopless graph on the vertex set $\{1,\dots,n\}$ (we use $\{0,\dots,n-1\}$
in the code). A \emph{perfect matching} is a set of edges covering every vertex
exactly once. An \emph{orientation} $D$ directs every edge, and its skew adjacency
matrix is $A_D=(a_{ij})$, where $a_{ij}=1$ if $i\to j$, $a_{ij}=-1$ if $j\to i$, and
$a_{ij}=0$ otherwise. Let $M=\{t_1h_1,\dots,t_mh_m\}$ be a perfect matching
($n=2m$), with each edge written tail-first ($t_k\to h_k$). Its \emph{term} is
\[
  \operatorname{sgn}_D(M)=\operatorname{sgn}\begin{pmatrix}1&2&\cdots&2m-1&2m\\
  t_1&h_1&\cdots&t_m&h_m\end{pmatrix}=(-1)^{\operatorname{inv}(t_1h_1\cdots t_mh_m)},
\]
where $\operatorname{inv}$ counts inversions. This does not depend on the order in
which the edges are listed: moving a block of two entries past another entry is an
even permutation. The Pfaffian is
\[
\operatorname{Pf}(A_D)=\sum_{\pi}\operatorname{sgn}(\pi)\prod_{k}a_{\pi(2k-1)\pi(2k)}
=\sum_{M}\operatorname{sgn}_D(M).
\]
Here $\pi$ runs over the permutations with $\pi(1)<\pi(3)<\dots$ and
$\pi(2k-1)<\pi(2k)$. A term is non-zero exactly when the pairs form a perfect
matching $M$. Writing each pair tail-first changes the sign of the permutation
and of the factor $a$ together, so the term equals $\operatorname{sgn}_D(M)$.

\begin{definition}
$D$ is a \emph{Pfaffian orientation} if $|\operatorname{Pf}(A_D)|$ equals the number
$\Phi(G)$ of perfect matchings. $G$ is \emph{Pfaffian} if it has a Pfaffian orientation.
\end{definition}

\begin{lemma}\label{lem:equiv}
$|\operatorname{Pf}(A_D)|=\Phi(G)$ if and only if all perfect matchings have the same sign
$\operatorname{sgn}_D(M)$.
\end{lemma}
\begin{proof}
Suppose $p$ of the $\Phi$ terms are $+1$ and $q$ are $-1$. Then
$|\operatorname{Pf}|=|p-q|$, and this equals $p+q$ if and only if $pq=0$.
\end{proof}

\section{The parity obstruction}
\begin{lemma}\label{lem:flip}
Reversing one edge $e$ of $D$ changes $\operatorname{sgn}_D(M)$ for exactly the perfect
matchings $M\ni e$.
\end{lemma}
\begin{proof}
If $e\notin M$, the word $t_1h_1\cdots t_mh_m$ is unchanged. If $e=t_kh_k\in M$, the
two adjacent entries $t_k,h_k$ are swapped. They are distinct because $G$ is
loopless, so the swap is a transposition and the inversion number changes by
exactly $1$.
\end{proof}

\begin{theorem}[Parity obstruction]\label{thm:obstruction}
Let $G$ be loopless, and suppose that
\begin{enumerate}
\item[(a)] $\Phi(G)$ is even;
\item[(b)] every edge lies in an even number of perfect matchings;
\item[(c)] for some orientation $D_0$, the number of perfect matchings $M$ with
$\operatorname{sgn}_{D_0}(M)=-1$ is odd.
\end{enumerate}
Then $G$ is not Pfaffian.
\end{theorem}
\begin{proof}
Let $\nu(D)$ be the number of negative terms. Any orientation $D$ is obtained from
$D_0$ by reversing edges one at a time. By Lemma~\ref{lem:flip}, reversing $e$
changes $\nu$ by an amount congruent modulo $2$ to the number of perfect matchings
containing $e$, which is even by (b). Hence $\nu(D)\equiv\nu(D_0)\equiv1\pmod2$ for
every $D$. If $D$ were Pfaffian, Lemma~\ref{lem:equiv} would give $\nu(D)\in\{0,\Phi\}$,
which is even by (a). This is a contradiction.
\end{proof}
Equivalently, by (b) the product $\prod_M\operatorname{sgn}_D(M)=(-1)^{\nu(D)}$ does not
depend on $D$. It is $-1$, while a Pfaffian orientation would make it
$(\pm1)^{\Phi}=1$. This is the classical reason why $K_{3,3}$ is not Pfaffian
(Little~\cite{little}; see also \cite{lp,rst}).

\section[K\_\{3,3\}]{\texorpdfstring{$K_{3,3}$}{K33}}
Let the vertices be $0,1,2$ and $3,4,5$, and let $D_0$ orient every edge $a\to b$ with
$a\in\{0,1,2\}$. The perfect matchings are $M_\sigma=\{0\,\sigma(0),\,1\,\sigma(1),\,2\,\sigma(2)\}$
for the six bijections $\sigma:\{0,1,2\}\to\{3,4,5\}$. In the word
$0\,\sigma(0)\,1\,\sigma(1)\,2\,\sigma(2)$ the entries $\sigma(i)\ge3$ form $2+1=3$
inversions with the later entries $1,2$, plus the inversions of $\sigma$ itself, so
$\operatorname{sgn}_{D_0}(M_\sigma)=-\operatorname{sgn}(\sigma)$:
\begin{center}
\begin{tabular}{lccc}
\toprule
matching & mask (Lean) & $\operatorname{sgn}(\sigma)$ & $\operatorname{sgn}_{D_0}$\\
\midrule
$\{05,14,23\}$ & 84 & $-1$ & $+1$\\
$\{04,15,23\}$ & 98 & $+1$ & $-1$\\
$\{05,13,24\}$ & 140 & $+1$ & $-1$\\
$\{03,15,24\}$ & 161 & $-1$ & $+1$\\
$\{04,13,25\}$ & 266 & $-1$ & $+1$\\
$\{03,14,25\}$ & 273 & $+1$ & $-1$\\
\bottomrule
\end{tabular}
\end{center}
(Bit $i$ of a mask selects edge number $i$ in the order $03,04,05,13,14,15,23,24,25$.)
Thus $\Phi=6$; each edge $ab$ lies in exactly $2$ matchings (the $2$ bijections with
$\sigma(a)=b$); and $\nu(D_0)=3$ is odd. By Theorem~\ref{thm:obstruction}:

\begin{proposition}\label{prop:k33}
$K_{3,3}$ is not Pfaffian. Concretely, $\operatorname{Pf}(A_{D_0})=0$, and every one of
the $2^9=512$ orientations has $|\operatorname{Pf}(A_D)|\le 4<6$.
\end{proposition}
(Indeed $\nu(D)$ is odd, so $|\operatorname{Pf}|=|6-2\nu|\in\{4,0\}$.)

\section[All even subdivisions of K\_\{3,3\}]{All even subdivisions of \texorpdfstring{$K_{3,3}$}{K33}}
Let $G$ have vertices $\{0,\dots,n-1\}$ and an edge $uv$. Let $G'$ be the graph
obtained by replacing $uv$ with the path $u\,a\,b\,v$, where $a=n$ and $b=n+1$ are
new vertices. Every even subdivision of $K_{3,3}$ arises from $K_{3,3}$ by finitely
many such steps, up to relabelling of the vertices.

\begin{lemma}\label{lem:sub}
If $G$ satisfies (a), (b) and (c) of Theorem~\ref{thm:obstruction}, so does $G'$, and
$\Phi(G')=\Phi(G)$.
\end{lemma}
\begin{proof}
In $G'$ the vertex $a$ is matched either to $b$ or to $u$; in the latter case $b$ is
matched to $v$. This gives a bijection $M\mapsto M'$ from the perfect matchings of $G$:
\[
M'=M\cup\{ab\}\ \text{if } uv\notin M,\qquad M'=(M\setminus\{uv\})\cup\{ua,bv\}\ \text{if } uv\in M.
\]
So $\Phi(G')=\Phi(G)$ and (a) is inherited. Let $c$ be the number of perfect
matchings of $G$ through $uv$. In $G'$ the edges $ua$ and $bv$ lie in $c$ perfect
matchings and $ab$ lies in $\Phi-c$. Both numbers are even, and the other edges keep
their counts, so (b) is inherited. For (c): since $G$ satisfies (b), the proof of
Theorem~\ref{thm:obstruction} shows that $\nu(D)$ is odd for \emph{every} orientation
$D$ of $G$, so we may take $D_0$ with $u\to v$. Orient
$u\to a\to b\to v$ and keep every other edge. If $uv\notin M$, the word of $M'$ is the word of
$M$ followed by $a\,b$. These are the two largest labels, in increasing order, so the
sign is unchanged. If $uv\in M$, list $uv$ last, so that the word of $M$ is $W\,u\,v$.
The word of $M'$ is $W\,u\,a\,b\,v$. It arises from $W\,u\,v\,a\,b$ (which has the sign of
$W\,u\,v$) by moving $v$ past $a$ and $b$, an even permutation. Hence
$\operatorname{sgn}_{D_0'}(M')=\operatorname{sgn}_{D_0}(M)$ for all $M$, and $\nu$ is
preserved.
\end{proof}

Relabelling the vertices by a permutation $\tau$ multiplies every term by
$\operatorname{sgn}(\tau)$. So conditions (a), (b) and (c) are invariant under
relabelling: when $\Phi$ is even, $\nu$ either stays the same or becomes $\Phi-\nu$.
Being Pfaffian is also invariant. Induction with Lemma~\ref{lem:sub}, starting from
Proposition~\ref{prop:k33}, gives:

\begin{theorem}\label{thm:main}
Every even subdivision of $K_{3,3}$ (including $K_{3,3}$ itself) satisfies the
hypotheses of Theorem~\ref{thm:obstruction} with $\Phi=6$. Hence no even subdivision
of $K_{3,3}$ is Pfaffian, and Clause (E) of Conjecture 00000001685 is false.
\end{theorem}

\begin{corollary}
Conjecture 00000001685 is false.
\end{corollary}

The smallest proper even subdivision is $S$: $K_{3,3}$ with edge $03$ replaced by
the path $0\,6\,7\,3$ (8 vertices, 11 edges). It is non-Pfaffian with
$\Phi(S)=6$. The masks of its perfect matchings, over the edge order
$04,05,13,14,15,23,24,25,06,67,73$, are $554,561,582,645,1360,1416$.

\section{M\"obius ladders}
The M\"obius ladder with $r$ rungs has vertices $0,\dots,2r-1$, the cycle edges
$i\,(i+1 \bmod 2r)$ and the rungs $i\,(i+r)$. For $r=3$ it is $K_{3,3}$, with colour
classes the even and the odd vertices. Its $6$ perfect matchings satisfy the
obstruction directly; this is proved in Lean as \texttt{M6\_not\_pfaffian}. An
exhaustive computation in \texttt{verify.py}, using the method of the next section,
gives the following:
\begin{center}
\begin{tabular}{lcccccc}
\toprule
rungs $r$ & 2 & 3 & 4 & 5 & 6 & 7\\
vertices & 4 & 6 & 8 & 10 & 12 & 14\\
Pfaffian? & yes & no & yes & no & yes & no\\
\bottomrule
\end{tabular}
\end{center}
The ladder with $5$ rungs has $13$ perfect matchings, and none of its $2^{15}$
orientations is Pfaffian. In the vertex-count convention, then, the M\"obius part
of the conjecture also fails (for example at $M_6$ and $M_{10}$). In the rung-count
convention we make no claim about it.

\section{Secondary: the second sentence}\label{sec:second}
Read the second sentence as ``every nonplanar Pfaffian graph of crossing number
$\le2$ is a M\"obius ladder or an even subdivision of $K_{3,3}$''. Under this
reading it is also false. Let $W$ be the Wagner graph $V_8$ (the M\"obius ladder
with $4$ rungs) with the edge $01$ replaced by a path $0\,8\,9\,1$. Then:
\begin{itemize}
\item $W$ is Pfaffian, with $\Phi(W)=7$. The proof of Lemma~\ref{lem:sub} shows that
the compatible orientation of a Pfaffian orientation of $V_8$ is Pfaffian. The
Lean theorem \texttt{M8\_pfaffian} exhibits a Pfaffian orientation of $V_8$, and
\texttt{verify.py} checks $W$ directly.
\item $W$ is nonplanar. $V_8-\{04\}$ is a subdivision of $K_{3,3}$ with classes
$\{1,3,6\}$ and $\{2,5,7\}$: the paths are $12$, $15$, $1\,0\,7$, $32$, $37$,
$3\,4\,5$, $65$, $67$, $62$.
\item $\operatorname{cr}(W)=1$. The Wagner graph has crossing number
$\operatorname{cr}(V_8)=1$: Guy and Harary~\cite{gh} show that every M\"obius ladder with at
least $3$ rungs has crossing number $1$. Crossing number is invariant under subdividing edges: a drawing
of $G$ is a drawing of any subdivision of $G$ with the same crossings, and
conversely, by suppressing the degree-$2$ vertices.
\item $W$ is not a M\"obius ladder, since it has vertices of degree $2$. It is not an
even subdivision of $K_{3,3}$, since $|E|-|V|=14-10=4$, whereas every subdivision of
$K_{3,3}$ has $|E|-|V|=3$.
\end{itemize}
This part is not formalized in Lean, apart from \texttt{M8\_pfaffian}. The main
disproof does not depend on it.

\section{Independent computational check (\texttt{verify.py})}
The script uses only the Python standard library and a different method from the
Lean proof. Since $\det A_D=\operatorname{Pf}(A_D)^2$ and $|\operatorname{Pf}|\le\Phi$, $G$ is
Pfaffian if and only if some orientation has $\det A_D=\Phi^2$. The script computes
exact integer determinants (Bareiss elimination) and enumerates perfect matchings
recursively. It checks:
\begin{itemize}
\item $K_{3,3}$ (all $512$ orientations), $S$ (all $2048$), and $K_{3,3}$ with two
edges subdivided twice (all $8192$): none is Pfaffian, and $\Phi=6$ in each case.
\item $K_{3,3}$ with every edge subdivided $2$ times, and with every edge subdivided
$4$ times. Switching at a vertex (reversing all its edges) multiplies the
Pfaffian by $-1$ and leaves $\det$ unchanged. So it suffices to check one
orientation per switching class: fix the edges of a spanning tree, which leaves
$2^{|E|-|V|+1}=16$ orientations. None is Pfaffian.
\item The hypotheses (a), (b), (c) of Theorem~\ref{thm:obstruction} for all $2^9$ even
subdivisions with $0$ or $2$ internal vertices per edge, and for three more with
up to $6$.
\item M\"obius ladders with $2$--$7$ rungs (the table above), with $5$ rungs also
exhaustively over $2^{15}$ orientations; and the graph $W$ of the previous
section.
\end{itemize}

\section{Lean formalization}
\texttt{lean4/} is a self-contained Lean~4.19.0 project (core library only, no
Mathlib). Everything is defined from scratch:
\begin{itemize}
\item \texttt{Graph} is a vertex count and an edge list (edge endpoints are assumed to
be $<n$, which holds for every graph used). \texttt{sel es M} is the edge
subset selected by the bit mask $M$. \texttt{isPM} says every vertex $v<n$ lies on
exactly one selected edge. \texttt{PMs G} lists all perfect matchings among the
$2^{|E|}$ masks, and \texttt{numPM G} is their number.
\item An orientation is a \texttt{List Bool}, one entry per edge (\texttt{true} keeps
$(a,b)$ as $a\to b$). \texttt{par} is the inversion parity of the word
$t_1h_1t_2h_2\cdots$, and \texttt{term} is $\pm1$ accordingly.
$\texttt{Pf G os}=\sum_{M}\texttt{term}$.
\item \texttt{IsPfaffian G}: some orientation of length $|E|$ has
$|\texttt{Pf}|=\texttt{numPM}$ (the conjecture's definition). \texttt{IsPfaffian'}:
all terms are equal (Kasteleyn's definition). \texttt{same\_sign\_pfaffian} shows
that the second implies the first.
\item \texttt{not\_pfaffian\_of\_invariant} is Theorem~\ref{thm:obstruction} for
arbitrary graphs. It is proved through \texttt{word\_perm}, \texttt{par\_flips}
(Lemma~\ref{lem:flip} in cumulative form: the parity changes by the number of
reversed matching edges), \texttt{sum\_map\_flips\_even} and \texttt{sum\_terms}.
\item Computations (kernel \texttt{decide}): \texttt{PMs\_K33}, \texttt{PMs\_S},
\texttt{PMs\_M6}, and the hypotheses (a)--(c). The results are
\texttt{K33\_not\_pfaffian}, \texttt{S\_not\_pfaffian} and \texttt{M6\_not\_pfaffian}.
\item \texttt{EvenSubdivisionK33} is the inductive family generated from $K_{3,3}$ by
\texttt{subdivide2}, which replaces an edge $uv$ by $u\,n\,(n{+}1)\,v$.
\texttt{S\_even\_subdivision} shows $S$ belongs to it and $S\ne K_{3,3}$.
\item Main theorems:
\texttt{conjecture\_00000001685\_false (Q)}:
$\neg\,\forall G,\ \texttt{EvenSubdivisionK33}\ G\to Q\,G\wedge\texttt{IsPfaffian}\ G$, for an
arbitrary extra property $Q$ such as ``nonplanar'' or ``minimal'';
the same with \texttt{IsPfaffian'}; \texttt{\_false\_proper} (members other than
$K_{3,3}$); \texttt{\_conjunction\_false} ($\neg((E)\wedge R)$ for any $R$); and
\texttt{mobius\_clause\_false\_vertex\_convention}; and for the obstruction reading,
\texttt{obstruction\_reading\_false} and
\texttt{obstruction\_reading\_false\_vertex\_convention}.
\item Non-vacuity: \texttt{K4\_pfaffian} ($K_4$, the ladder with $2$ rungs) and
\texttt{M8\_pfaffian} (the nonplanar Wagner graph) have explicit Pfaffian
orientations, so \texttt{IsPfaffian} is not empty. As a cross-check that
\texttt{Pf} is the matrix Pfaffian, \texttt{pfMat\_K33} and \texttt{pfMat\_M8} compare
it with the recursive row expansion
$\operatorname{Pf}(A)=\sum_j(-1)^{j}a_{1j}\operatorname{Pf}(A_{\hat1\hat j})$ of the skew
matrix.
\end{itemize}
\paragraph{Not in Lean.} Lean covers $K_{3,3}$ and $S$. These suffice for every
reading in which the clause applies to all even subdivisions (with or without
$K_{3,3}$ itself). The induction to all even subdivisions
(Lemma~\ref{lem:sub}) is proved on paper only, as is the invariance under
relabelling, and Section~\ref{sec:second} apart from \texttt{M8\_pfaffian}.
There is no
\texttt{sorry}, no \texttt{native\_decide}, and no added axiom.
\texttt{\#print axioms} shows at most \texttt{propext} and \texttt{Quot.sound}.

\section{Reproduction}
\begin{verbatim}
cd lean4 && lake build     # Lean 4.19.0, core only (about 1 minute)
cd .. && python3 verify.py # Python 3 standard library
pdflatex report.tex && pdflatex report.tex
\end{verbatim}

\begin{thebibliography}{9}
\bibitem{little} C.~H.~C.~Little, A characterization of convertible $(0,1)$-matrices,
\emph{J. Combin. Theory Ser. B} 18 (1975), 187--208.
\bibitem{gh} R.~K.~Guy and F.~Harary, On the M\"obius ladders, \emph{Canad. Math. Bull.}
10 (1967), 493--496.
\bibitem{lp} L.~Lov\'asz and M.~D.~Plummer, \emph{Matching Theory}, North-Holland, 1986.
\bibitem{rst} N.~Robertson, P.~D.~Seymour and R.~Thomas, Permanents, Pfaffian
orientations, and even directed circuits, \emph{Ann. of Math.} 150 (1999), 929--975.
\end{thebibliography}
\end{document}
